% !TEX root = ../../main.tex
%% Corrigé du défi du chapitre 15 — vérifié par verifications/defis_11_15.py.

\begin{blocLivre}{cor-15-25a}
\begin{corrige}{25}{L'intégrale de Gauss}
\demarche{Encadrer l'intégrande par deux fonctions dont l'intégrale est un
Wallis, puis laisser l'équivalent de Wallis faire converger les deux
bornes vers la même limite.}
\partieExo{A}{Existence et encadrement}
\begin{questions}
  \item $F'(A)=\mathrm{e}^{-A^{2}}>0$ : $F$ est croissante. Pour
    $t\geqslant1$, $t^{2}\geqslant t$ donc $\mathrm{e}^{-t^{2}}\leqslant\mathrm{e}^{-t}$,
    et pour $A\geqslant1$ :
    $F(A)\leqslant\int_0^{1}1\,\mathrm{d}t+\int_1^{A}\mathrm{e}^{-t}\,\mathrm{d}t
    \leqslant1+\frac1{\mathrm{e}}$. Croissante et majorée, $F$ admet une
    \res{\text{limite finie }G} en $+\infty$ (théorème de la limite
    monotone).
  \item $\varphi(u)=\mathrm{e}^{u}-1-u$ a pour dérivée $\mathrm{e}^{u}-1$,
    du signe de $u$ : minimum $\varphi(0)=0$, donc $\varphi\geqslant0$.
    Avec $u=-\frac{t^{2}}{n}$ : $0\leqslant1-\frac{t^{2}}{n}\leqslant\mathrm{e}^{-t^{2}/n}$
    pour $t\leqslant\sqrt n$, et l'on élève à la puissance $n$ (fonction
    croissante sur $\R_{+}$). Avec $u=\frac{t^{2}}{n}$ :
    $\bigl(1+\frac{t^{2}}{n}\bigr)^{n}\leqslant\mathrm{e}^{t^{2}}$, puis on
    passe à l'inverse.
\end{questions}
\partieExo{B}{Le calcul par les intégrales de Wallis}
\begin{questions}[resume]
  \item $t=\sqrt n\sin\theta$, $\theta\in\intFF{0}{\frac\pi2}$ :
    $\mathrm{d}t=\sqrt n\cos\theta\,\mathrm{d}\theta$ et
    $1-\frac{t^{2}}{n}=\cos^{2}\theta$, d'où
    $\displaystyle\int_0^{\sqrt n}\Bigl(1-\frac{t^{2}}{n}\Bigr)^{n}\mathrm{d}t
    =\sqrt n\int_0^{\pi/2}\cos^{2n+1}\theta\,\mathrm{d}\theta=\res{\sqrt n\,W_{2n+1}}$.
  \item $t=\sqrt n\tan\theta$, $\theta\in\intFF{0}{\frac\pi4}$ :
    $\mathrm{d}t=\frac{\sqrt n}{\cos^{2}\theta}\,\mathrm{d}\theta$ et
    $\bigl(1+\tan^{2}\theta\bigr)^{-n}=\cos^{2n}\theta$, d'où l'intégrale
    $\sqrt n\displaystyle\int_0^{\pi/4}\cos^{2n-2}\theta\,\mathrm{d}\theta
    \leqslant\res{\sqrt n\,W_{2n-2}}$ (intégrande positif, intervalle plus
    grand).
  \item En intégrant l'encadrement de la question 2 sur
    $\intFF{0}{\sqrt n}$ : $\sqrt n\,W_{2n+1}\leqslant F(\sqrt n)\leqslant\sqrt n\,W_{2n-2}$.
    Or
    $\sqrt n\,W_{2n+1}\sim\sqrt n\sqrt{\frac{\pi}{4n+2}}\to\frac{\sqrt\pi}{2}$
    et de même $\sqrt n\,W_{2n-2}\sim\sqrt n\sqrt{\frac{\pi}{4n-4}}\to\frac{\sqrt\pi}2$.
    Par encadrement $F(\sqrt n)\to\frac{\sqrt\pi}{2}$, et comme $F$ a pour
    limite $G$ : $\res{G=\dfrac{\sqrt\pi}{2}}\approx0{,}8862$.
\end{questions}
\end{corrige}
\end{blocLivre}

\begin{blocLivre}{cor-15-25b}
\begin{corrigeSuite}
\partieExo{C}{Applications}
\begin{questions}[start=6]
  \item Avec $u=\sqrt a\,t$ :
    $\displaystyle\int_0^{A}\mathrm{e}^{-at^{2}}\mathrm{d}t
    =\frac{1}{\sqrt a}\int_0^{\sqrt aA}\mathrm{e}^{-u^{2}}\mathrm{d}u
    \xrightarrow[A\to+\infty]{}\res{\frac{\sqrt\pi}{2\sqrt a}}$.
  \item Par parité puis question 6 avec $a=\frac12$ :
    $\displaystyle\int_{-A}^{A}\mathrm{e}^{-t^{2}/2}\,\mathrm{d}t
    =2\int_0^{A}\mathrm{e}^{-t^{2}/2}\,\mathrm{d}t\to2\cdot\frac{\sqrt\pi}{2}\sqrt2=\sqrt{2\pi}$,
    d'où la limite \res{1}. Numériquement, l'intégrale entre $-1{,}96$ et
    $1{,}96$ vaut déjà environ $0{,}95$ : la règle des « $95\,\%$ » des
    statisticiens.
  \item Le même $\sqrt\pi$ vient chaque fois des intégrales de Wallis, donc
    de $\int_0^{\pi/2}$ : dans Stirling ($n!\sim\sqrt{2\pi n}\,(n/\mathrm{e})^{n}$),
    dans le binôme central ($\binom{2n}{n}\sim4^{n}/\sqrt{\pi n}$, d'où la
    probabilité de retour $P(S_{2n}=0)\sim1/\sqrt{\pi n}$ de la marche
    aléatoire), et ici. Ce n'est pas une coïncidence : la loi binomiale
    renormalisée converge vers la loi normale, dont la densité est
    $\frac{1}{\sqrt{2\pi}}\mathrm{e}^{-t^{2}/2}$ — trois chapitres, un seul
    phénomène.
\end{questions}
\end{corrigeSuite}
\end{blocLivre}
